% Преамбула отчёта (ГОСТ 7.32).
% Класс G7-32 уже загружает: fontspec, polyglossia (русский/английский),
% amsmath, caption (настройки ГОСТ), geometry (поля ГОСТ), titlesec, tikz,
% multirow, etoolbox. Повторно их не подключаем.

% Моноширинный шрифт (fontspec загружается классом). Liberation Mono входит в
% образ сборки (пакет fonts-liberation) и содержит кириллицу.
\setmonofont{Liberation Mono}

% ГОСТ 7.0.11 требует правое поле 15 мм (класс по умолчанию ставит 10 мм).
\geometry{right=15mm}

% Более мягкая вёрстка абзацев: разрешаем растяжение межсловных интервалов,
% чтобы длинные идентификаторы кода (\texttt) переносились на новую строку,
% а не выходили за правое поле.
\sloppy
\setlength{\emergencystretch}{3em}

% Математика
\usepackage{amsfonts}
\usepackage{amssymb}

% Графика
\usepackage{graphicx}
\usepackage{float}

% Схемы (архитектура вычислительного слоя, конвейер симуляции)
\usepackage{tikz}
\usetikzlibrary{arrows.meta, positioning}
% Кинематические схемы (глава о генерации кинематики из STEP)
\usepackage{kinematikz}

% Таблицы
\usepackage{booktabs}
\usepackage{array}
\usepackage{longtable}

% Цвет
\usepackage{xcolor}

% Листинги исходного кода
\usepackage{listings}
\definecolor{lstkw}{rgb}{0.0,0.0,0.5}
\definecolor{lstcomment}{rgb}{0.30,0.45,0.30}
\definecolor{lststring}{rgb}{0.55,0.10,0.10}
\lstset{
  basicstyle=\footnotesize\ttfamily,
  keywordstyle=\color{lstkw}\bfseries,
  commentstyle=\color{lstcomment}\itshape,
  stringstyle=\color{lststring},
  breaklines=true,
  frame=single,
  numbers=left,
  numberstyle=\tiny,
  showstringspaces=false,
  tabsize=2,
  captionpos=b,
  columns=fullflexible,
  keepspaces=true,
}
% Подпись листингов на русском
\renewcommand{\lstlistingname}{Листинг}

% --- Псевдокод ---------------------------------------------------------------
% Алгоритмы оформляются псевдокодом в обрамлённом блоке; нумерация общая с
% листингами. Использование:
%   \begin{pseudocode}{<подпись>}{<метка>}
%     строки псевдокода, отступ уровня n --- \pind{n}, перенос строки --- \\
%   \end{pseudocode}
\newcommand{\pind}[1]{\hspace*{#1\dimexpr1.6em\relax}}
\newsavebox{\pseudobox}
\newenvironment{pseudocode}[2]{%
  \par\medskip\noindent
  \begin{minipage}{\linewidth}%
  \def\pseudocaption{#1}\def\pseudolabel{#2}%
  \setlength{\fboxsep}{8pt}%
  \begin{lrbox}{\pseudobox}%
  \begin{minipage}{\dimexpr\linewidth-2\fboxsep-2\fboxrule\relax}%
  \small\setlength{\parindent}{0pt}\raggedright
}{%
  \end{minipage}\end{lrbox}%
  \noindent\fbox{\usebox{\pseudobox}}\par\smallskip
  \captionof{lstlisting}{\pseudocaption}\label{\pseudolabel}%
  \end{minipage}\par\medskip
}

% Гиперссылки (подключать последними)
\usepackage{hyperref}
\hypersetup{
  colorlinks=true,
  linkcolor=black,
  citecolor=black,
  urlcolor=blue,
  pdftitle={Промежуточный научно-технический отчёт по этапу 1},
  pdfauthor={farfield},
}

% --- Заглушка под скриншот -------------------------------------------------
% Использование:
%   \screenshot{<доля \textwidth>}{<подпись рисунка>}{<что заснять>}{<метка>}
% Рисует обрамлённую область с пояснением; нумеруется как обычный рисунок и
% попадает в список иллюстраций. Когда изображение готово, заглушку заменяют
% на \includegraphics из каталога figures/.
\newcommand{\screenshot}[4]{%
  \begin{figure}[ht]\centering
    \setlength{\fboxsep}{0pt}%
    \fbox{\begin{minipage}[c][0.40\textwidth][c]{#1\textwidth}%
      \centering
      \itshape\small[\,ЗАГЛУШКА ДЛЯ СКРИНШОТА\,]\par\medskip
      \normalfont\small #3\par\medskip
      \footnotesize\ttfamily figures/#4.png
    \end{minipage}}%
    \caption{#2}\label{fig:#4}
  \end{figure}%
}
